% Part: lambda-calculus
% Chapter: introduction
% Section: syntax

\documentclass[../../../include/open-logic-section]{subfiles}

\begin{document}

\olfileid{lam}{int}{syn}
\olsection{د لامبډا حساب نحو}

لومړے د متغيرونو يوه لړۍ $x$، $y$، $z$،~\dots او ځينې ثابتې نښې
$a$، $b$، $c$،~\dots اخلو۔ د ترمونو سټ په استقرايي ډول داسې
تعريفېږي:
\begin{enumerate}
\item هر متغير يو ترم دے۔
\item هر ثابت يو ترم دے۔
\item که $M$ او $N$ ترمونه وي، $(MN)$ هم ترم دے۔
\item که $M$ ترم او $x$ متغير وي، $(\lambd[x][M])$ هم ترم دے۔
\end{enumerate}
د $(MN)$ په بڼه ترمونه \emph{تطبيقونه} او د
$(\lambd[x][M])$ په بڼه ترمونه \emph{تجريدونه} بلل کېږي۔

هغه نظام چې هېڅ ثابت نلري، \emph{خالص} لامبډا حساب بلل کېږي۔ موږ به
زياتره په خالص $\lambd$-حساب کښې کار کوو؛ نو ټول واړه لاتيني توري
به متغيرونه ښيي۔ د لويو لاتيني تورو ($M$، $N$ او داسې نورو) په وسيله
د $\lambd$-حساب ترمونه ښيو۔

د نښو څو دودونه به خپل کړو:

\begin{conv}
\begin{enumerate}
\item چې قوسونه پرېښودل شي، تطبيق له کيڼ لوري ښي لوري ته ترسره کېږي۔
  مثلاً، که $M$، $N$، $P$ او $Q$ ترمونه وي، نو $MNPQ$ د
  $(((MN)P)Q)$ لنډه بڼه ده۔
\item همداسې، چې قوسونه پرېښودل شي، لامبډا تجريد ته تر وسه پراخه
  ساحه ورکوو۔ مثلاً، $\lambd[x][MNP]$ د
  $(\lambd[x][((MN)P)])$ په معنا لولو۔
\item يوه لامبډا د څو متغيرونو د تجريد دپاره کارېدے شي۔ مثلاً،
  $\lambd[xyz][M]$ د
  $\lambd[x][\lambd[y][\lambd[z][M]]]$ لنډه بڼه ده۔
\end{enumerate}
\end{conv}

د بېلګې په توګه،
\[
\lambd[xy][xxyx \lambd[z][xz]]
\]
د دې لنډه بڼه ده:
\[
\lambd[x][\lambd[y][((((xx)y)x)(\lambd[z][(xz)]))]].
\]
دا دودونه په ياد کړئ۔ په لومړي سر کښې به درته ډېر ستړي کوونکي وي،
خو ورسره به عادت شئ؛ څه موده وروسته به تر بې‌شمېره قوسونو کم ستړي
کوونکي درته ښکاري۔

دوه ترمونه چې يوازې د تړلو متغيرونو په نومونو کښې توپير ولري،
$\alpha$-معادل بلل کېږي؛ مثلاً $\lambd[x][x]$ او
$\lambd[y][y]$۔ آسانه به وي چې دوئ ``هماغه'' ترم وګڼو؛ په بله
وينا، چې کله $M$ او $N$ يو شان بولو، د تړلو متغيرونو تر نوم بدلولو
پورې يو شانوالی هم مراد لرو۔ هغه متغيرونه چې د يوې $\lambd$ په ساحه
کښې وي، ``تړلي''، او نور ``آزاد'' بلل کېږي۔ په مخکينۍ بېلګه کښې
آزاد متغير نشته؛ خو په
\[
(\lambd[z][yz])x
\]
کښې $y$ او $x$ آزاد دي او $z$ تړلے دے۔

\end{document}
